\documentclass[a4paper,11pt]{article}
\usepackage[utf8]{inputenc}
\usepackage[T1]{fontenc}
\usepackage[french]{babel}
\usepackage[left=2cm,right=2cm,top=2cm,bottom=2cm]{geometry}
\usepackage{xcolor}
\usepackage{hyperref}
\usepackage{fontawesome5}
\usepackage{parskip}

% --- Couleurs EasyMile ---
\definecolor{primary}{RGB}{0, 82, 147}
\hypersetup{colorlinks=true, urlcolor=primary}

\begin{document}

% --- EXPÉDITEUR ---
\noindent
\begin{minipage}[t]{0.5\textwidth}
    \textbf{Loïc Aron MBASSI EWOLO}\\
    Élève-Ingénieur Bac+5 -- Double diplôme ENSEA\\[3pt]
    \faMapMarker\ Cergy, France\\
    \faPhone\ (+33) 7 59 52 33 98\\
    \faEnvelope\ \href{mailto:loicaron.mbassiewolo@ensea.fr}{loicaron.mbassiewolo@ensea.fr}
\end{minipage}
\hfill
% --- DATE ---
\begin{minipage}[t]{0.4\textwidth}
    \raggedleft
    Cergy, le 10 février 2026
\end{minipage}

\vspace{1.2cm}

% --- DESTINATAIRE ---
\hfill
\begin{minipage}[t]{0.5\textwidth}
    \raggedleft
    \textbf{EasyMile}\\
    Équipe Perception\\
    Toulouse, France
\end{minipage}

\vspace{1cm}

\noindent\textbf{Objet :} Candidature au stage Deep Learning on accumulated LiDAR PointCloud

\vspace{0.8cm}

Madame, Monsieur,

Le traitement de nuages de points Lidar 3D denses pour la segmentation sémantique et la création de HD Maps constitue un défi technique que je travaille actuellement dans le cadre du challenge CoHoMa (robotique autonome militaire). Cette expérience, combinée à mon stage de recherche au MILA sur les réseaux de neurones, m'a préparé aux missions proposées dans votre équipe Perception.

Dans le cadre du challenge CoHoMa à l'ENSEA, je traite des données issues d'un Lidar 3D VLP16 pour la cartographie d'environnements non structurés. Ce travail m'a permis de manipuler des nuages de points volumineux, d'implémenter des algorithmes de SLAM, et de déployer des modèles de détection (YOLOv10) sur Nvidia Jetson avec optimisation GPU. L'approche Superpoint Transformer que vous mentionnez pour la segmentation 3D efficace correspond directement à la problématique de passage à l'échelle que j'ai rencontrée.

Mon stage au MILA (Institut Québécois d'IA) m'a apporté une rigueur dans l'implémentation de pipelines d'entraînement PyTorch et l'analyse de la convergence des réseaux. Cette méthodologie de recherche sera utile pour évaluer les architectures candidates (Superpoint Transformer, EZ-SP) et adapter les données EasyMile au format d'entraînement requis.

Je suis disponible dès avril 2026 pour une durée de 5 à 6 mois. La perspective de contribuer à l'amélioration du stack de perception offline, notamment pour l'auto-annotation d'obstacles statiques et la génération de HD Maps, représente une opportunité de mettre mes compétences au service de la mobilité autonome.

Je reste à votre disposition pour un entretien afin de discuter de ma candidature.

Je vous prie d'agréer, Madame, Monsieur, l'expression de mes salutations distinguées.

\vspace{0.8cm}

\textbf{Loïc Aron MBASSI EWOLO}\\
\textit{Élève-Ingénieur ENSEA / Polytechnique Yaoundé}

\end{document}
